% =========================================================
%  config/commands.tex
%  Perintah khusus buku ajar.
%  Tujuan: mencegah definisi berulang yang tersebar di chapters/.
%  Perintah \CPL, \CPMK, \SUBCPMK, \code, \istilah, \headrow,
%  \Utama, \Pendukung, \TidakAda, \tabsource diambil dari
%  rps/config/commands.tex dan rps/config/styles.tex agar
%  penulisan buku ajar dan RPS konsisten.
% =========================================================

% --- Kode capaian pembelajaran --------------------------
\newcommand{\CPL}[1]{\textbf{CPL#1}}
\newcommand{\CPMK}[1]{\textbf{CPMK#1}}
\newcommand{\SUBCPMK}[1]{\textbf{Sub-CPMK#1}}

% --- Istilah teknis -------------------------------------
% \istilah  : istilah asing (miring)
% \code     : potongan kode, nama berkas, atau nama fungsi di dalam teks
%
% \code memakai \detokenize sehingga garis bawah dan karakter khusus lain
% boleh ditulis APA ADANYA tanpa perlu di-escape:
%     benar : \code{pygame.key.get_pressed}   \code{code/python/uji_gerak.py}
%     salah : \code{pygame.key.get\_pressed}  (akan tercetak "\_")
\newcommand{\code}[1]{\texttt{\detokenize{#1}}}
\newcommand{\istilah}[1]{\textit{#1}}

% --- Penanda baris header tabel -------------------------
\newcommand{\headrow}{}

% --- Penanda status hubungan (matriks CPL--CPMK) --------
\newcommand{\Utama}{\textbf{Utama}}
\newcommand{\Pendukung}{Pendukung}
\newcommand{\TidakAda}{---}

% --- Penanda indikator pencapaian -----------------------
% Dipakai pada tabel Sub-CPMK: \ind{a}, \ind{b}, \ind{c}
\newcommand{\ind}[1]{\textbf{(#1)}}

% --- Keterangan sumber gambar ---------------------------
% Contoh: \sumber{[R12] hal. 24}
\newcommand{\sumber}[1]{%
  \par\vspace{2pt}{\footnotesize\itshape Sumber: #1\par}%
}

% --- Akar direktori yang bergantung pada cara kompilasi ----
% Ketika sebuah bab dikompilasi mandiri (misalnya latexmk -pdf chapter.tex
% di dalam chapters/05), direktori kerjanya adalah direktori bab sehingga
% akar buku berada dua tingkat di atas. Ketika seluruh buku dikompilasi,
% direktori kerjanya adalah akar buku.
%
% Kondisi \ifSubfilesClassLoaded dievaluasi SEKALI di sini, yaitu ketika
% berkas ini dibaca sebagai bagian dari preamble.
%
% PENTING: \ifSubfilesClassLoaded BUKAN kondisional gaya \if...\else...\fi.
% Perintah ini menerima DUA argumen dalam kurung kurawal, seperti
% \@ifundefined: \ifSubfilesClassLoaded{<bila mandiri>}{<bila seluruh buku>}.
% Menulisnya sebagai \if...\else...\fi menyebabkan galat "Extra \else" dan
% menghasilkan nama berkas yang kacau.
%
% Kondisi juga tidak boleh diletakkan di dalam argumen nama berkas
% (\includegraphics, \lstinputlisting) karena argumen tersebut diproses
% ulang oleh \edef.
\ifSubfilesClassLoaded{%
  \newcommand{\coderoot}{../../code/}%
  \newcommand{\assetroot}{../../assets/}%
}{%
  \newcommand{\coderoot}{code/}%
  \newcommand{\assetroot}{assets/}%
}

% Dipakai sebagai: \lstinputlisting{\codefile{python/ch04_game_loop.py}}
\newcommand{\codefile}[1]{\coderoot#1}

% Dipakai sebagai: \includegraphics{\assetfile{images/tangkapan.png}}
\newcommand{\assetfile}[1]{\assetroot#1}

% --- Rujukan ke bab lain ---------------------------------
% \refbab{chap:game-loop}{Bab tentang Game Loop}
% Mencetak nomor bab bila label tersedia (kompilasi seluruh buku atau
% label diimpor melalui paket xr-hyper); bila tidak, mencetak teks
% cadangan sehingga bab tetap dapat dikompilasi mandiri (Arsitektur §40).
\makeatletter
\newcommand{\refbab}[2]{%
  \@ifundefined{r@#1}{#2}{Bab~\ref{#1}}%
}
\makeatother

% --- Penanda bagian yang belum selesai (hanya untuk penyusunan) --
% Tidak boleh dipakai pada versi terbit.
\newcommand{\todo}[1]{%
  \begin{tcolorbox}[
    breakable, sharp corners, boxrule=0.5pt,
    colframe=ubRed, colback=ubRedBg, coltitle=white,
    fonttitle=\bfseries\small, title={Belum selesai}]
  \footnotesize #1
  \end{tcolorbox}%
}
